%-------------------------------------------------------------------------------
% midi_learn
%-------------------------------------------------------------------------------
%
% \file        midi_learn.tex
% \library     Documents
% \author      Chris Ahlstrom
% \date        2026-06-22
% \update      2026-07-09
% \version     $Revision$
% \license     $XPC_GPL_LICENSE$
%
%     This section discusses the details of the MIDI Learn functionality.
%
%-------------------------------------------------------------------------------

\section{MIDI Learn}
\label{sec:midi_learn}

   The new (as of version 0.99.26) MIDI Learn feature of
   \textsl{Seq66} provides rudimentary support for a
   mapping of MIDI events to many of the controls of the
   \textsl{Seq66} user-interface.

   Hitherto, adding these mappings required looking up
   a MIDI controllers output events in its programmer's
   manual, and then hand-editing them into a
   \textsl{Seq66} 'ctrl' file.
   See \sectionref{subsec:configuration_ctrl}.
   See \sectionref{sec:launchpad_mini}.

   In the \textsl{Seq66} implementation, one must first bring up
   the "learn manager" dialog, select the type of events to
   be learned, and what kind of action is to be added via
   the learning process.
   The process is described in detail later.

   There are some limitations and quirks in the learn features:

   \begin{itemize}
      \item There is currently no way to make the learnable
         controls obvious. If an action is learnable, it
         shows up in the learn manager.
      \item Ignoring the live-grid slots, there are currently only
         80 or so user-interface actions that can be mapped.
         Not everything clickable can be mapped,
         and currently there is no way know without trying them,
         though the most common visible actions are mappable.
         This could be expanded somewhat, pending user requests.
      \item There are no "continuous controls" (such as a volume
         or pan knob) in \textsl{Seq66}.
         Hence no direct mappings of CC events.
         However, they can still be used to trigger
         an action.
      \item Making corrections, especially for loops and mute-groups,
         is clumsy.
      \item There is no way to use MIDI Learn to set up output
         control events to affect lights on the MIDI controller.
   \end{itemize}

   We are still working on coordinating and improving the MIDI
   Learn process.

   Prior to using MIDI learn, the MIDI controls from the 'ctrl' are
   loaded at start-up, and copied into the "performer" as the main
   control set, if set as active in the 'rc'.
   Remember that a control container has three sections for input
   controls.

   \begin{itemize}
      \item \textbf{[loop-control]}
      \item \textbf{[mute-group-control]}
      \item \textbf{[automation-control]}
   \end{itemize}

   We'll call these "Categories".

   In each section, present is the name of a hot-key, plus
   three "stanzas" that specify the inverse flag, the event status and
   value, plus a range of acceptable values.
   The sections are:

   \begin{itemize}
      \item \textbf{[toggle]}
      \item \textbf{[on]}
      \item \textbf{[off]}
   \end{itemize}

   When the \textbf{Learn Manager} is brought up, the main controls
   are copied into a separate ("original") container to support the
   "Reset" operation, and then copied to another ("current")
   container for modification.
   When events are added, they are added to the "current" container,
   replacing any event that matches the current combination of
   category Loops/Mutes/Automation and Toggle/On/Off.

   The following dialog can be brought to the
   screen by the \textbf{Learn} menu (see\sectionref{subsec:menu_learn}.)

\begin{figure}[H]
   \centering 
   \includegraphics[scale=0.65]{main-window/learn-manager-2.png}
   \caption{Seq66 Learn Manager}
   \label{fig:menu_learn_manager}
\end{figure}

   While this dialog is open, the MIDI Learn functionality is available,
   and it goes away when the dialog closes.

   \setcounter{ItemCounter}{0}      % Reset the ItemCounter for this list.

   \itempar{Learn Categories}{learn!categories}
   This radio-group determines into which section the current MIDI control
   click is sent. Whenever one of the learn categories is selected,
   that starts the process of mapping MIDI events to actions.
   The following briefly describes how each "category" works.

   \begin{itemize}
      \item \textbf{Loop}.
         Starting at pattern (loop) 0 or just after the last pattern,
         the user simply presses a button on the MIDI control grid,
         then keeps pressing additional buttons until up to 32 buttons
         have been mapped. This operation is simple and fast.
      \item \textbf{Mute Groups}.
         Starting at mute-group 0 or just after the last mute-group,
         the user simply presses a button on the MIDI control grid,
         then keeps pressing additional buttons until up to 32 buttons
         have been mapped. This operation is simple and fast.
      \item \textbf{Automation}.
         In this mode, the user clicks the desired GUI element,
         verifies that it's effect is shown in the Learn Manager,
         and then clicks the desired MIDI controller button.
         This sequence is required for each GUI element to be learned.
   \end{itemize}

   \itempar{Active Events}{learn!active events}
   These boxes show the number of \textsl{active} MIDI control input events.
   These are events that have a non-zero MIDI status byte.

   \itempar{Learn Action}{learn!action}
   This item selects which action the MIDI event will trigger:

   \begin{itemize}
      \item \textbf{Toggle}.
         The default is "Toggle", which is the most common trigger needed.
         Most hot-keys act as toggles, as an example.
         The MIDI control press is logged as a "toggle".
         This should be the action to use for the "Loop" and
         "Mute Groups" mappings.
      \item \textbf{On}.
         If selected, the MIDI control press is allocated as an "on" action.
      \item \textbf{Off}.
         If selected, the MIDI control press is allocated as an "off" action.
      \item \textbf{Hold}.
         A "hold" selection logs the MIDI control press as an "on"
         action, and logs the MIDI control release as an "off" action.
         Useful for maintaining a certain state only when the MIDI
         controller button is held down.
   \end{itemize}

   \itempar{Learn Values}{learn!values}

   \begin{itemize}
      \item \textbf{Inverse}.
         "Inverse" reverses the action of the control.
         Almost never used, as far as we know.
      \item \textbf{d1min}.
         Determines the minimum d1 (second data byte) value to
         trigger the action.
      \item \textbf{d1max}.
         Determines the maximum d1 (second data byte) value to
         trigger the action.
   \end{itemize}

   \textbf{d1man} and \textbf{d1max}, determine
   the allowable values of, for example, note velocity.
   In most cases, we want the largest range, 0-127.

   \itempar{Learn Save}{learn!save}
   When clicked, this copies the currently-entered events into
   the application's main control-in container, and also
   replaces the "original" container to support a possible
   reset operation.
   In addition, the 'ctrl' file is marked for saving.

   \itempar{Learn Reset}{learn!Reset}
   When clicked, this copies the original or last-saved events into
   the current container.

   \itempar{Learn Clear All}{learn!clear all}
   When clicked, all events are cleared from the current container.
   This includes the loop, mute-group, and automation events.

   \itempar{Learn Clear}{learn!clear selected}
   When clicked, all events in the selected category
   (Loop, Mute Groups, or Automation) \textsl{only}
   are cleared from the current container.
   The text of the button indicates what will be cleared.

   \itempar{Learn OK}{learn!OK}
   When clicked, acts like the \textbf{Save} button, then closes
   the Learn Manager.

   \itempar{Learn Cancel}{learn!cancel}
   When clicked, simply closes the Learn Manager, with no changes
   since the last save.

   The \textbf{Loops} and \textbf{Mutes} control recordings work similarly;
   they don't require clicking on the user-interface.
   The \textbf{Automation} control recording requires a little more
   work.

\subsection{Learn / Loops}
\label{subsec:midi_learn_loops}

   In the user interface, the muting and unmuting of a particular pattern in
   the live grid is handled using either a hot-key or by clicking on a
   button in the grid.
   The Learn-Loops process allows easy assignment of a MIDI event
   to each grid slot.

   When activated, this process works by auto-incrementing the loop number
   from 0 to 31 as each MIDI controller button is pressed.
   Obviously, this works best on a grid controller.

   \begin{enumber}
      \item Bring up the Learn Manager and make sure the \textbf{Loops}
         radio button is active.
      \item Make sure the Loops \textbf{Active Events} count is zero.
         If not, click the \textbf{Clear} button, which clears
         \textsl{all} controls, or click the \textbf{Clear loop}.
         Beware!
      \item Verify that the \textbf{Next Control} box shows "Loop 0".
      \item Leave the \textbf{Action} set to \textsl{Toggle},
         the \textbf{Inverse} unchecked, \textbf{d1min} set to 0,
         and \textbf{d1max} set to 127.
         This ensures good recording and acceptance of all note velocities,
         for example.
      \item Click the first button on the controller to log that event
         to toggle pattern 0.
      \item Click the remaining buttons all the way up to pattern 31.
         One can stop early, if desired.
         Verify that the \textbf{Active Events} count are
         as expected.
      \item If some mistakes were made, the only recourse is to
         click the \textbf{Reset} button and start over.
      \item Click the \textbf{Save} button.
         This copies the current events to the main control-in set.
   \end{enumber}

\subsection{Learn / Mute Groups}
\label{subsec:midi_mute_groups}

   The process here is similar to the previous section, except that the
   controls involve mute-group numbers, not pattern/loop numbers.

   When activated, this process works by auto-incrementing the
   mute-group number
   from 0 to 31 as each MIDI controller button is pressed.
   Obviously, this works best on a grid controller.

   \begin{enumber}
      \item Bring up the Learn Manager and make sure the \textbf{Mutes}
         radio button is active.
      \item Make sure the Mutes \textbf{Active Events} count is zero.
         If not, click the \textbf{Clear} button, which clears
         \textsl{all} controls, or click the \textbf{Clear mutegroup}.
         Beware!
      \item Verify that the \textbf{Next Control} box shows
         "Mute Group 0".
      \item Leave the \textbf{Action} set to \textsl{Toggle},
         the \textbf{Inverse} unchecked, \textbf{d1min} set to 0,
         and \textbf{d1max} set to 127.
         This ensures good recording and acceptance of all note velocities,
         for example.
      \item Click the first button on the controller to log that event
         to toggle mute-group 0.
      \item Click the remaining buttons all the way up to mute-group 31.
         One can stop early, if desired.
         Verify that the \textbf{Active Events} count are
         as expected.
      \item If some mistakes were made, the only recourse is to
         click the \textbf{Reset} button and start over.
      \item Click the \textbf{Save} button.
         This copies the current events to the main control-in set.
   \end{enumber}

\subsection{Learn / Automation}
\label{subsec:midi_learn_automation}

   Automation controls work differently than the loop and mute-group grid
   handling described above.
   The desired hot-key or user-interface element must be clicked first,
   followed by clicking the desired MIDI control button.

   \begin{enumber}
      \item Bring up the Learn Manager and make sure the
         \textbf{Automation} radio button is active.
      \item Make sure the Automation \textbf{Active Events} count is zero.
         If not, click the \textbf{Clear} button, which clears
         \textsl{all} controls, or click the \textbf{Clear automation}.
         Beware!
      \item Verify that the \textbf{Next Control} box shows
         "Automation" followed by the name of the desired control.
      \item Leave the \textbf{Action} set to \textsl{Toggle},
         the \textbf{Inverse} unchecked, \textbf{d1min} set to 0,
         and \textbf{d1max} set to 127.
         This ensures good recording and acceptance of all note velocities,
         for example.
      \item Click the first button on the controller to log that event
         to toggle mute-group 0.
      \item Click the remaining buttons all the way up to mute-group 31.
         One can stop early, if desired.
         Verify that the \textbf{Active Events} count are
         as expected.
      \item If some mistakes were made, the only recourse is to
         click the \textbf{Reset} button and start over.
      \item Click the \textbf{Save} button.
         This copies the current events to the main control-in set.
   \end{enumber}

\subsection{Learn / Automation Cheat Diagrams}
\label{subsec:midi_learn_automation_cheat_diagrams}

   This section shows the \textsl{Seq66} user-interface elements that
   can be learned, and the seq66::automation::slot enumeration values
   for each. When the element is clicked, that enumeration value
   is shown in the learn manager.

\begin{figure}[H]
   \centering 
   \includegraphics[scale=0.50]{midi-learning/main-window.png}
   \caption{Main Window Automation Elements}
   \label{fig:main_window_automation_elements}
\end{figure}

\begin{figure}[H]
   \centering 
   \includegraphics[scale=0.50]{midi-learning/pattern-window.png}
   \caption{Pattern Window Automation Elements}
   \label{fig:pattern_window_automation_elements}
\end{figure}

\begin{figure}[H]
   \centering 
   \includegraphics[scale=0.50]{midi-learning/song-window.png}
   \caption{Song Window Automation Elements}
   \label{fig:song_window_automation_elements}
\end{figure}

%-------------------------------------------------------------------------------
% vim: ts=3 sw=3 et ft=tex
%-------------------------------------------------------------------------------
